\labsection{3}{Reshape time-aware data and engineer defensible features}{sec:lab03}

\begin{labproject}{Lab 3 wrangling and feature engineering}
\textbf{Dataset.} \texttt{bike\_sharing\_dataset\_hour.csv}.\\
\textbf{Goal.} Transform a large hourly table into features and summaries that answer a specific demand question.

\textbf{Procedure.}
\begin{enumerate}
  \item Parse the date field and document the row grain: one row represents one observed hour.
  \item Create readable day-type and hour labels without destroying the original coded variables.
  \item Engineer peak-period categories and cyclical sine/cosine encodings for hour-of-day.
  \item Aggregate mean demand by hour and working-day status, reducing more than 17,000 rows to an interpretable profile.
  \item Compare casual and registered demand across time periods and explain why aggregating these columns separately matters.
  \item Check the engineered cyclical features mathematically and visually.
\end{enumerate}

\textbf{Interpretation checkpoint.} Feature engineering changes the representation, not the underlying evidence. Every engineered variable should have a clear definition that another analyst can reproduce.

\lstinputlisting[style=dsr]{all_experiments/lab03_wrangling.R}
\end{labproject}